% Part: proof-theory
% Chapter: sequent-calculus
% Section: quantifiers

\documentclass[../../../include/open-logic-section]{subfiles}

\begin{document}

\olfileid{pt}{seq}{qua}

\olsection{નિયમિત સાબિતીઓ અને પ્રતિસ્થાપન}

પરિમાણકના નિયમોમાં કેટલીક જટિલતાઓ ઊભી થાય છે,
કારણ કે \LeftR{\lexists} અને~\RightR{\lforall} પર
``આઇગનચલ-શરતો'' લાગુ પડે છે.
આવા અનુમાનોમાં !!{constant}~$a$નો ક્યારેય પુનઃઉપયોગ ન થાય
તે સુનિશ્ચિત કરીને મોટાભાગની જટિલતાઓ ઉકેલી શકાય.
આ વ્યાખ્યા રજૂ કરીએ.

\begin{defn}
\Log{G1c} કે \Log{G3c}માં !!^a{proof}~$\pi$ \emph{નિયમિત} છે, જો
\begin{enumerate}
  \item કોઈ આઇગનચલ~$a$ એક કરતાં વધુ
  \LeftR{\lexists} કે~\RightR{\lforall} અનુમાનનું આઇગનચલ ન હોય, અને
  \item $a$ એ \LeftR{\lexists} કે~\RightR{\lforall}નું આઇગનચલ હોય,
  તો તે $\pi$માં ફક્ત એ અનુમાનની ઉપર જ આવે.
\end{enumerate}
\end{defn}

\begin{lem}\ollabel{lem:subst-regular}
$\pi$ નિયમિત હોય અને તેનો અંત $\Gamma \Sequent \Delta$ હોય,
$c$ એ~$\pi$માં આઇગનચલ તરીકે ન વપરાતો !!a{constant} હોય,
અને $t$~એ~$\pi$નું કોઈ આઇગનચલ ન ધરાવતું બંધ પદ હોય,
તો~$\pi$માં $c$ની દરેક આવૃત્તિને બદલે~$t$ મૂકવાથી મળતું
$\Subst{\pi}{t}{c}$ નિયમિત !!{proof} છે.
$\Subst{\pi}{t}{c}$નો અંતિમ સિક્વન્ટ
$\Subst{\Gamma}{t}{c} \Sequent \Subst{\Delta}{t}{c}$ છે,
જ્યાં $\Subst{\Gamma}{t}{c} =
\Setabs{!A(t)}{!A(c) \in \Gamma}$.
સંદર્ભો બહુગણો હોવાથી આ રૂપાંતરમાં દરેક આવૃત્તિ જળવાય છે.
\sourcecorrection{OLMD-329}{મૂળ પ્રતિસ્થાપિત પદ માટે ફક્ત આઇગનચલો ન ધરાવવાની શરત આપે છે. આ અધ્યાયના નિયમ-વિભાગમાં પરિમાણકનો દાખલ બંધ પદ હોવાનો છે; ખુલ્લું પદ મૂકવાથી એ શરત અથવા પરિમાણકની અંદર ચલ-બંધન બદલાઈ શકે. તેથી અહીં બંધ પદની શરત સ્પષ્ટ કરી છે.}
\end{lem}

\begin{proof}
  $\pheight{\pi}$ પર આગમનથી સાબિતી કરીએ.
  $\pheight{\pi} = 0$ હોય, તો ફક્ત સ્વયંસિદ્ધ છે.
  ત્યારે $\Subst{\pi}{t}{c}$ પણ સ્વયંસિદ્ધ છે,
  કારણ કે તેનું કોઈ !!{formula}~$!A(c)$ એ $!A(t)$ બને છે.

  $\pheight{\pi} > 0$ હોય,
  તો~$\pi$ના છેલ્લા અનુમાન પ્રમાણે કિસ્સા જુદા પાડીએ.
  બંધારણીય નિયમોના (\Log{G1c}માં) અને વિધાનતર્કીય સંયોજકોના
  તાર્કિક નિયમોના કિસ્સા સરળ છે અને સ્વાધ્યાય તરીકે છોડીએ.
  $\pi$નો અંત પરિમાણકના અનુમાનથી થાય તે કિસ્સા વિચારીએ.
  \begin{enumerate}
    \item છેલ્લું અનુમાન $\RightR{\lforall}$ હોય,
    તો $\pi$નું સ્વરૂપ આ છે:
\[
    \AxiomC{}
    \RightLabel{$\pi'$}
    \Deduce$\Gamma(c) \fCenter \Delta'(c), !A(a,c)$
    \RightLabel{\RightR{\lforall}}
    \UnaryInf$\Gamma(c) \fCenter \Delta'(c), \lforall[x][!A(x,c)]$
    \DisplayProof
    \]
    આગમન-ધારણાથી $\Subst{\pi'}{t}{c}$ એ
    $\Gamma(t) \fCenter \Delta'(t), !A(a,t)$ની
    નિયમિત !!{proof} છે.
    $a$ એ~$\pi$નું આઇગનચલ હોવાથી $a$~એ~$t$માં આવતું નથી.
    એટલે $\Subst{\pi}{t}{c}$નું છેલ્લું અનુમાન,
\[
    \AxiomC{}
    \RightLabel{$\Subst{\pi'}{t}{c}$}
    \Deduce$\Gamma(t) \fCenter \Delta'(t), !A(a,t)$
    \RightLabel{\RightR{\lforall}}
    \UnaryInf$\Gamma(t) \fCenter \Delta'(t), \lforall[x][!A(x,t)]$
    \DisplayProof
    \]
    યોગ્ય છે: નિષ્કર્ષમાં~$a$ આવતું નથી.
    \item \Log{G3c}માં છેલ્લું અનુમાન $\LeftR{\lforall}$ હોય,
    તો $\pi$નું સ્વરૂપ આ છે:
\[
    \AxiomC{}
    \RightLabel{$\pi'$}
    \Deduce$!A(s(c),c), \lforall[x][!A(x,c)], \Gamma(c) \fCenter \Delta'(c)$
    \RightLabel{\LeftR{\lforall}}
    \UnaryInf$\lforall[x][!A(x,c)], \Gamma(c) \fCenter \Delta'(c)$
    \DisplayProof
    \]
    આગમન-ધારણાથી $\Subst{\pi'}{t}{c}$ એ
    $!A(s(t),t), \lforall[x][!A(x,t)], \Gamma(t) \fCenter \Delta'(t)$ની
    નિયમિત !!{proof} છે.
    $\Subst{\pi}{t}{c}$નું છેલ્લું અનુમાન,
\[
    \AxiomC{}
    \RightLabel{$\Subst{\pi'}{t}{c}$}
    \Deduce$!A(s(t),t), \lforall[x][!A(x,t)], \Gamma(t) \fCenter \Delta'(t)$
    \RightLabel{\LeftR{\lforall}}
    \UnaryInf$\lforall[x][!A(x,t)], \Gamma(t) \fCenter \Delta'(t)$
    \DisplayProof
    \]
    યોગ્ય છે. \Log{G1c}માં~\LeftR{\lforall}નો કિસ્સો સમાન,
    પરંતુ થોડો સરળ છે.
    \sourcecorrection{OLMD-330}{મૂળના આ સર્વપરિમાણક-ડાબા કિસ્સાનાં બંને વૃક્ષોમાં જમણા નિયમનું લેબલ છે. મુખ્ય સૂત્ર તથા આધારમાં જાળવેલી તેની આવૃત્તિ મુજબ બંને લેબલ ડાબા કર્યા છે.}
    \item છેલ્લું અનુમાન \RightR{\lexists} કે \LeftR{\lexists} હોય: સ્વાધ્યાય.
  \end{enumerate}
  દરેક કિસ્સામાં નિયમિત !!{proof}ના અંતે સિક્વન્ટ ઉમેર્યો છે
  (બે આધારવાળા નિયમોમાં બે નિયમિત !!{proof}sના અંતે).
  નવું સિક્વન્ટ $\pi$ના અંતિમ સિક્વન્ટથી ફક્ત એટલું જુદું છે
  કે $c$ને બદલે પદ~$t$ મૂક્યું છે.
  વળી $\pi$ અને $\Subst{\pi}{t}{c}$ના આઇગનચલો એ જ છે.
  $t$માં~$\pi$નું કોઈ આઇગનચલ નથી એમ ધાર્યું હોવાથી,
  નવા અંતિમ સિક્વન્ટમાં $\Subst{\pi}{t}{c}$નું કોઈ આઇગનચલ નથી.
  એટલે $\Subst{\pi}{t}{c}$ નિયમિત છે.
\end{proof}

\begin{prop}
$\pi$ એ \Log{G1c} કે \Log{G3c}ની !!a{proof} હોય,
તો એ જ સંરચના (ખાસ કરીને એ જ !!{height}) ધરાવતી
નિયમિત !!{proof}~$\pi'$ છે,
જે $\pi$થી ફક્ત આઇગનચલોનાં પ્રતિસ્થાપનથી જુદી પડે છે.
\end{prop}

\begin{proof}
ફક્ત આ સાબિતી પૂરતું, $\pi$ના આઇગનચલ-અનુમાનને \emph{સ્વચ્છ} કહીએ,
જો તેનું આઇગનચલ બીજા અનુમાનનું આઇગનચલ ન હોય
અને~$\pi$માં તે અનુમાનની તરત પહેલાંની ઉપ-!!{proof}
સિવાય ક્યાંય ન આવે; નહીં તો \emph{અસ્વચ્છ} કહીએ.
$\pi$માં અસ્વચ્છ આઇગનચલ-અનુમાનોની સંખ્યા $n$ લો.
$n$ પર આગમનથી બતાવીએ કે $\pi$~ને
જરૂરી નિયમિત !!{proof}~$\pi'$માં ફેરવી શકાય.
$n=0$ હોય, તો~$\pi$નાં બધાં આઇગનચલ-અનુમાન સ્વચ્છ છે;
એટલે $\pi$ નિયમિત છે અને કશું સાબિત કરવાનું નથી.
નહિતર સૌથી ઉપરનું અસ્વચ્છ આઇગનચલ-અનુમાન પસંદ કરો,
દાખલા તરીકે \RightR{\lforall}.
તે અનુમાનથી પૂરી થતી $\pi$ની ઉપ-!!{proof}~$\pi_1$નું સ્વરૂપ આ છે:
\[
    \AxiomC{}
    \RightLabel{$\pi_2$}
    \Deduce$\Gamma \fCenter \Delta, !A(c)$
    \RightLabel{\RightR{\lforall}}
    \UnaryInf$\Gamma \fCenter \Delta, \lforall[x][!A(x)]$
    \DisplayProof
    \]
$c$ આઇગનચલ હોવાથી~$\Gamma$, $\Delta$
કે $\lforall[x][!A(x)]$માં આવતું નથી તે નોંધો.
સાબિતી $\pi_2$ નિયમિત છે:
પસંદ કરેલા અસ્વચ્છ અનુમાનની ઉપરના આઇગનચલ-અનુમાન સ્વચ્છ છે.
પસંદ કરેલું આઇગનચલ તેમાં બીજા અનુમાનનું આઇગનચલ નથી.
\sourcecorrection{OLMD-331}{મૂળ સૌથી ઉપરનું આઇગનચલ-અનુમાન પસંદ કરે છે, પરંતુ તે સ્વચ્છ પણ હોઈ શકે, એટલે અસ્વચ્છ અનુમાનોની સંખ્યા ઘટવાની ખાતરી નથી. સૌથી ઉપરનું અસ્વચ્છ અનુમાન પસંદ કર્યું છે. તેની ઉપર સ્વચ્છ અનુમાન હોઈ શકે; તેથી ઉપસાબિતી નિયમિત છે તે કહ્યું છે, એકેય આઇગનચલ-અનુમાન નથી એવો મૂળનો દાવો કર્યો નથી.}
$\pi$માં ન આવતો !!a{constant} $a$ લો.
\olref{lem:subst-regular} મુજબ $\Subst{\pi_2}{a}{c}$
એ $\Gamma \fCenter \Delta, !A(a)$ની નિયમિત સાબિતી છે;
તેના પર~$\RightR{\lforall}$ લગાવી શકીએ.
$\pi$માં $\pi_1$ને બદલે સાબિતી $\pi_1'$,
\[
    \AxiomC{}
    \RightLabel{$\Subst{\pi_2}{a}{c}$}
    \Deduce$\Gamma \fCenter \Delta, !A(a)$
    \RightLabel{\RightR{\lforall}}
    \UnaryInf$\Gamma \fCenter \Delta, \lforall[x][!A(x)]$
    \DisplayProof
    \]
મૂકીએ, તો અસ્વચ્છ આઇગનચલ-અનુમાનને બદલે સ્વચ્છ અનુમાન મૂક્યું છે;
ફક્ત પસંદ કરેલા આઇગનચલ $c$ને બદલે~$a$ મૂકવાનો તફાવત છે.
મળેલી સાબિતીમાં વધુમાં વધુ $n-1$ અસ્વચ્છ આઇગનચલ-અનુમાનો છે,
અને પરિણામ આગમન-ધારણાથી અનુસરે છે.
\sourcecorrection{OLMD-332}{મૂળ નવી સંખ્યા બરાબર એકથી ઘટે છે કહે છે. બીજા અનુમાનોની એક જ આઇગનચલના પુનઃઉપયોગની ખામી પણ આ બદલાવથી દૂર થઈ શકે; તેથી ઓછામાં ઓછા એક અસ્વચ્છ અનુમાનનો ઘટાડો, એટલે વધુમાં વધુ જૂની સંખ્યાથી એક ઓછાં, એવો પૂરતો દાવો કર્યો છે.}
\end{proof}

હમણાં સાબિત કરેલા પ્રતિજ્ઞાપનને આગળ સ્પષ્ટ રીતે ટાંકશું નહીં,
પરંતુ વિચારેલી બધી !!{proof}s નિયમિત છે એમ માનીશું:
તે ન હોય, તો આઇગનચલો બદલીને હંમેશાં નિયમિત બનાવી શકીએ.

\end{document}
